\documentclass[a4paper,11pt]{report}
% define the title
\usepackage{graphicx}
\usepackage[brazil]{babel}
\usepackage[utf8]{inputenc}		%Usar acentos diretos do teclado
\usepackage[T1]{fontenc}		%Padrão para imagens
\usepackage{subfigure}
\usepackage{rotating}
\usepackage{longtable}
\usepackage{amsmath}
\author{Pedro Gushiken}
\date{4 de Abril de 2015}
\title{Controle Adaptativo de Velocidade de uma Máquina CC em vazio pela regra do MIT e pelo m\'{e}todo da sensibilidade}
\begin{document}
% generates the title
\maketitle

\tableofcontents
\chapter{Introdução}
\paragraph{Modelo da Máquina}
O modelo da máquina CC é descrito pelas seguintes equações de estado:
\begin{subequations}

$
X=
\begin{bmatrix}
I_{a} \\
W_{r}
\end{bmatrix}
$

~

$
U=
\begin{bmatrix}
V_{a}\\
T_{m}
\end{bmatrix}
$

~

$
\begin{bmatrix}
\dot{I}_{a} \\
\dot{W}_{r}
\end{bmatrix}
=
\begin{bmatrix}
-\frac{R_{a}}{L_{a}}-\frac{K_{e}\lambda_{e}}{L_{a}} \\
\frac{K_{e}\lambda_{e}}{J_{m}}-\frac{F_{m}}{J_{m}}
\end{bmatrix}
\begin{bmatrix}
I_{a}\\
W_{r}
\end{bmatrix}
+
\begin{bmatrix}
1/L_{a} & 0 \\
0 & -1/J_{m}
\end{bmatrix}
\begin{bmatrix}
V_{a}\\
T_{m}
\end{bmatrix}
$
\end{subequations}

Escolhemos controlar a máquina apenas a partir da tensão de armadura $V_{a}$ e a partir deste ponto projetamos um controlador tendo em mente o cancelamento de polos. E usamos os seguintes parâmetros para a máquina:
K_{e}=476.16


\paragraph{Modelo do Controlador com parâmetros conhecidos}
\chapter{Implementação do Controle Adaptativo através da regra do MIT}

\chapter{Implementação do Controle Adaptativo através do método da sensibilidade}

\chapter{Convergência dos parâmetros do controlador}
\paragraph{Para uma referência constante}
\paragraph{Para uma referência persistentemente excitante}
\chapter{Conclusão}
% insert the table of contents
\end{document}